\documentclass[12pt,a4paper]{article}
\usepackage{amsmath,amsthm,amssymb,amsfonts}
\usepackage[T1]{fontenc}
\usepackage[utf8]{inputenc}
\usepackage{hyperref}
\usepackage{geometry}
\geometry{margin=2.5cm}

\newtheorem{theorem}{Theorem}[section]
\newtheorem{lemma}[theorem]{Lemma}
\newtheorem{proposition}[theorem]{Proposition}
\newtheorem{corollary}[theorem]{Corollary}
\newtheorem{definition}[theorem]{Definition}
\newtheorem{remark}[theorem]{Remark}
\newtheorem{conjecture}[theorem]{Conjecture}

\begin{document}

\title{Tao's Probabilistic Approach to the Collatz Conjecture:\\
Almost All Orbits Attain Almost Bounded Values}
\author{Michael Fuhrmann}
\date{March 2026}
\maketitle

\begin{abstract}
In 2022, Terence Tao published a landmark result on the Collatz conjecture:
for any function $f\colon \mathbb{N} \to \mathbb{R}_{>0}$ with $f(n) \to \infty$,
almost all positive integers $n$ (in the sense of logarithmic density) eventually
have a Collatz iterate below $f(n)$. This paper presents Tao's approach in detail,
covering the probabilistic random-walk model for the Syracuse function, the Fourier-analytic
tools used to handle correlations between steps, the logarithmic density framework, and
the key estimate that shows divergent Collatz orbits have logarithmic density zero.
We also discuss why Tao's method falls short of proving the full conjecture and what
would be needed to bridge the remaining gap.
\end{abstract}

\tableofcontents

\newpage

\section{Introduction}

Terence Tao's 2022 paper ``Almost all Collatz orbits attain almost bounded values''
\cite{Tao2022}, published in \emph{Forum of Mathematics, Sigma}, represents the most
significant advance on the Collatz conjecture in decades. It does not prove the conjecture,
but it establishes the strongest known unconditional result in its favour.

The key insight is probabilistic: the Syracuse function $S(n) = (3n+1)/2^{\nu_2(3n+1)}$
acts like a ``random'' walk with a strong downward drift, and Tao quantifies this drift
precisely using Fourier analysis over the 2-adic integers.

\subsection{Main result}

\begin{theorem}[Tao 2022]\label{thm:tao_main}
For any function $f\colon \mathbb{N} \to \mathbb{R}_{>0}$ with $f(n) \to \infty$ as $n \to \infty$,
\[
  \mathrm{logdens}\bigl(\{ n \in \mathbb{N} : \exists k \geq 0,\; T^k(n) \leq f(n) \}\bigr) = 1.
\]
Equivalently, for any $\varepsilon > 0$ and any $f \to \infty$, the set of $n$ for which
no iterate of $T$ falls below $f(n)$ has logarithmic density 0.
\end{theorem}

\section{The Syracuse Random Walk Model}

\begin{definition}[Logarithmic density]\label{def:logdens}
The \emph{logarithmic density} of a set $A \subseteq \mathbb{N}$ is
\[
  \mathrm{logdens}(A) \;=\; \lim_{N \to \infty} \frac{1}{\log N} \sum_{\substack{n=1 \\ n \in A}}^{N} \frac{1}{n},
\]
if this limit exists. Logarithmic density is dominated by the natural density but is
more sensitive to the distribution of small primes.
\end{definition}

\begin{definition}[Syracuse function]\label{def:syracuse}
For an odd positive integer $m$, the \emph{Syracuse function} is
\[
  S(m) \;=\; \frac{3m+1}{2^{\nu_2(3m+1)}},
\]
where $\nu_2(3m+1)$ denotes the 2-adic valuation (the number of times 2 divides $3m+1$).
The value $E(m) := \nu_2(3m+1)$ is the number of ``even steps'' consumed in passing from
$m$ to $S(m)$ in the standard Collatz iteration.
\end{definition}

\begin{proposition}[Distribution of $E(m)$]\label{prop:dist_E}
For a random odd positive integer $m$, the random variable $E(m) = \nu_2(3m+1)$ follows
approximately a geometric distribution with parameter $1/2$:
\[
  \Pr[E(m) = k] \;\approx\; 2^{-k} \quad \text{for } k \geq 1.
\]
More precisely, among odd integers $m \in [1, N]$, the fraction with $E(m) = k$ is
$2^{-k}(1 + O(N^{-1/2}))$ as $N \to \infty$.
\end{proposition}

\begin{proof}
For an odd integer $m$, we have $m \equiv 1$ or $3 \pmod{4}$.
If $m \equiv 3 \pmod{4}$: $3m+1 \equiv 10 \equiv 2 \pmod{4}$, so $\nu_2(3m+1) = 1$, i.e.\ $E(m)=1$.
If $m \equiv 1 \pmod{4}$: $3m+1 \equiv 4 \equiv 0 \pmod{4}$, so $\nu_2(3m+1) \geq 2$, i.e.\ $E(m)\geq 2$.
Thus $E(m) = 1$ iff $m \equiv 3 \pmod{4}$: probability $1/2$ among odd integers.
More generally, $E(m) \geq k$ iff $2^k \mid 3m+1$, which (since $\gcd(3,2)=1$) is equivalent to
$m \equiv r_k \pmod{2^k}$ for a uniquely determined odd residue $r_k$.
This happens with probability $2^{-(k-1)}$ among odd integers.
Thus $\Pr[E(m) = k] = \Pr[E(m) \geq k] - \Pr[E(m) \geq k+1] = 2^{-(k-1)} - 2^{-k} = 2^{-k}$.
\end{proof}

\begin{definition}[Log-scale random walk]\label{def:rw}
To model the ``size'' of the Syracuse iterates, take natural logarithms ($\log = \ln$).
For odd $m$,
\[
  \ln S(m) \;=\; \ln m + \ln 3 - E(m) \ln 2 + O(1/m).
\]
The increments $X_k = \ln 3 - E_k \ln 2$ (where $E_k = E(S^{k-1}(m))$) form the
individual steps of the random walk. The expected value is
\[
  \mathbb{E}[X_k] \;=\; \ln 3 - \ln 2 \cdot \mathbb{E}[E_k] = \ln 3 - 2\ln 2 = \ln(3/4) < 0.
\]
Since $\ln(3/4) \approx -0.288 < 0$, the walk has a strict negative drift.
An equivalent formulation in base-2 logarithms gives $\log_2(3/4) = \log_2 3 - 2 \approx -0.415$
per Syracuse step; both formulations express the same downward drift.
\end{definition}

\begin{theorem}[Negative drift implies convergence a.s.]\label{thm:drift}
Let $(X_k)_{k \geq 1}$ be i.i.d.\ random variables with $\mathbb{E}[X_1] = \ln(3/4) < 0$.
Then the random walk $W_n = \sum_{k=1}^n X_k$ satisfies $W_n \to -\infty$ almost surely.
In particular, for any fixed threshold $L$, the walk eventually falls below $L$ with
probability 1.
\end{theorem}

\begin{proof}
This is an immediate consequence of the strong law of large numbers: $W_n/n \to \mathbb{E}[X_1] = \ln(3/4) < 0$
almost surely. Since $\mathbb{E}[X_1] < 0$, we have $W_n \to -\infty$ a.s.
\end{proof}

\begin{remark}\label{rem:not_independent}
The crucial difficulty is that the steps $X_k = \ln 3 - E(S^{k-1}(m)) \ln 2$ are
\emph{not} independent: the distribution of $E(S^k(m))$ depends on the value of
$S^k(m)$, which depends on the entire history of the orbit. Tao's main contribution
is to handle these dependencies via Fourier analysis.
\end{remark}

\section{Fourier Analysis over the 2-adic Integers}

\begin{definition}[Characters of $\mathbb{Z}/2^k\mathbb{Z}$]\label{def:characters}
For $k \geq 1$ and $\xi \in \mathbb{Z}/2^k\mathbb{Z}$, the \emph{character}
$\chi_{\xi,k}\colon \mathbb{Z}/2^k\mathbb{Z} \to \mathbb{C}^*$ is
\[
  \chi_{\xi,k}(n) \;=\; e^{2\pi i \xi n / 2^k}.
\]
The 2-adic Fourier transform of a function $f\colon \mathbb{Z}_2 \to \mathbb{C}$
on residues mod $2^k$ is $\hat{f}(\xi) = \sum_{n \bmod 2^k} f(n) \chi_{\xi,k}(n)$.
\end{definition}

\begin{definition}[Syracuse exponential sum]\label{def:exp_sum}
For $\xi \in \mathbb{R}/\mathbb{Z}$ and $N$ an odd positive integer, the
\emph{Syracuse exponential sum} is
\[
  F_k(\xi) \;=\; \mathbb{E}_{m \bmod 2^k,\, m \text{ odd}}\bigl[ e^{2\pi i \xi \log_2 S(m)} \bigr],
\]
where the expectation is over odd residues mod $2^k$.
\end{definition}

\begin{lemma}[Tao's key estimate]\label{lem:key}
For all $\xi \in \mathbb{R}/\mathbb{Z}$ with $\|\xi\|_{\mathbb{R}/\mathbb{Z}} \geq \delta > 0$,
there exists $c(\delta) > 0$ such that
\[
  |F_k(\xi)| \;\leq\; e^{-c(\delta) k}
\]
for all sufficiently large $k$. This exponential decay in $k$ is the heart of Tao's proof.
\end{lemma}

\begin{proof}[Proof sketch (informal — not a rigorous proof)]
% BUG-B7-P31-EN/DE-001 behoben: Klar als informal deklariert; ξ-Bedingung korrigiert
\emph{Warning: The following is a heuristic sketch of the i.i.d.\ approximation
only.  A rigorous proof requires Tao's decoupling argument; see~\cite{Tao2022}.
The step from sum to product is not immediate and the independence assumption
is not satisfied exactly.}

The exponential sum $F_k(\xi)$ factors \emph{approximately} as a product over
``independent'' steps. If the $E_j$ were exactly i.i.d.\ geometric(1/2), then
$F_k(\xi) = \left(\sum_{j=1}^\infty 2^{-j} e^{2\pi i \xi (\log_2 3 - j)}\right)^k$,
and the inner sum is a geometric series in $e^{-2\pi i \xi}$:
\[
  F_k(\xi) = \left(\frac{e^{2\pi i \xi \log_2 3}}{2 - e^{-2\pi i \xi}}\right)^k.
\]
For $\|\xi\|_{\mathbb{R}/\mathbb{Z}} > 0$ (i.e.\ $\xi \not\in \mathbb{Z}$, which is
guaranteed by the lemma hypothesis $\|\xi\|_{\mathbb{R}/\mathbb{Z}} \geq \delta > 0$),
the denominator $2 - e^{-2\pi i \xi}$ has modulus $< 2$, so the modulus of the
inner factor is $< 1$, giving exponential decay in $k$.
Tao handles the non-independence via a decoupling argument using the Fourier
structure of $\mathbb{Z}_2$, which provides the necessary rigour for the
transition from the i.i.d.\ model to the actual (dependent) Syracuse steps.
\end{proof}

\section{The Logarithmic Density Argument}

\begin{definition}[Divergent set]\label{def:divergent}
Define the \emph{divergent set} as
\[
  \mathcal{D} \;=\; \bigl\{ n \in \mathbb{N} : T^k(n) \geq f(n) \text{ for all } k \geq 0 \bigr\},
\]
i.e.\ the set of integers whose Collatz orbit never descends below $f(n)$.
Tao's theorem asserts $\mathrm{logdens}(\mathcal{D}) = 0$.
\end{definition}

\begin{theorem}[Logarithmic density of divergent orbits]\label{thm:logdens_zero}
For any $f\colon \mathbb{N} \to \mathbb{R}_{>0}$ with $f(n) \to \infty$, the set
$\mathcal{D}$ has logarithmic density zero.
\end{theorem}

\begin{proof}[Proof outline (following Tao \cite{Tao2022})]
The proof proceeds in three main steps.

\textbf{Step 1: Reduction to the Syracuse map.}
One works with the compressed Syracuse function $S$ on odd integers. The Collatz iterate
$T^k(n)$ is at most $f(n)$ iff $S^j(n) \leq C \cdot f(n)$ for some $j$ and constant $C$.
Thus it suffices to show that the logarithmic density of odd $n$ with all Syracuse iterates
above $f(n)$ is zero.

\textbf{Step 2: Discretization via 2-adic residues.}
For each scale $N$, one studies the orbit of $n$ modulo $2^k$ for $k = O(\log \log N)$.
The Syracuse function acts on residues mod $2^k$, and the Fourier estimates of
Lemma~\ref{lem:key} show that correlations between orbit steps decay exponentially in $k$.

\textbf{Step 3: Exponential sum estimate.}
One writes the indicator of ``orbit stays above $f(n)$'' as a Fourier integral, applies
Lemma~\ref{lem:key} to bound the Fourier coefficients, and integrates. The bound
$|F_k(\xi)| \leq e^{-c(\delta)k}$ implies that the logarithmic density of integers with
non-descending orbits (after $k$ Syracuse steps) is at most $C e^{-c(\delta)k}$, which
tends to 0 as $k \to \infty$.
\end{proof}

\begin{corollary}[Collatz orbits are almost bounded]\label{cor:almost_bounded}
For any $\varepsilon > 0$, the set of $n \in \mathbb{N}$ for which $T^k(n) > n^\varepsilon$
for all $k \geq 0$ has logarithmic density zero. In particular, almost all (in logarithmic density)
Collatz orbits contain arbitrarily small values relative to the starting point.
\end{corollary}

\section{The Role of Ergodic Theory}

\begin{definition}[Syracuse random walk as ergodic process]\label{def:ergodic_rw}
The sequence of log-ratios $(\log S^k(n) / \log n)_{k \geq 0}$ can be viewed as the
trajectory of a point in a measure-preserving dynamical system.
\textbf{Important caveat:} The Syracuse map $S$ is \emph{not} measure-preserving
with respect to the Haar measure on $\mathbb{Z}_2^{\text{odd}}$.
To apply the Birkhoff ergodic theorem, one must work with a separately constructed
$S$-invariant measure $\nu$ on $\mathbb{Z}_2^{\text{odd}}$ (as done rigorously in
\cite{Tao2022} and in the companion paper on ergodic methods \cite{Kontorovich2010}).
Under such an invariant measure $\nu$, the function
$g(x) = \log_2 3 - E(x)$ (where $E(x) = \nu_2(3x+1)$) is the ``step size''
function with $\nu$-expectation
\[
  \int g\, d\nu = \log_2 3 - \mathbb{E}_\nu[E].
\]
If $\nu$ is chosen so that $E$ has a geometric distribution with parameter $1/2$
(which holds in the i.i.d.\ approximation), then $\mathbb{E}_\nu[E] = 2$ and
\[
  \frac{1}{k} \sum_{j=0}^{k-1} g(S^j(x)) \;\longrightarrow\; \log_2 3 - 2 = \log_2(3/4) < 0
\]
for $\nu$-almost every $x$ by the Birkhoff ergodic theorem.
\end{definition}

\begin{theorem}[Ergodic convergence of Syracuse orbits]\label{thm:ergodic_conv}
Let $\nu$ be an $S$-invariant ergodic measure on $\mathbb{Z}_2^{\text{odd}}$ under which
$E(x) = \nu_2(3x+1)$ has expectation $2$ (as in the i.i.d.\ geometric$(1/2)$ model).
Then for $\nu$-almost every $x$:
\[
  \lim_{k \to \infty} \frac{\log S^k(x)}{k} \;=\; \log(3/4) < 0.
\]
This means that the logarithm of the orbit decreases linearly, i.e.\ the orbit decreases
\emph{geometrically fast} in the absolute value sense.

\textbf{Note:} The Haar measure on $\mathbb{Z}_2^{\text{odd}}$ is \emph{not} preserved
by $S$, so the Birkhoff theorem cannot be applied directly with Haar measure.
The statement above uses the invariant measure $\nu$; see \cite{Kontorovich2010}
for an explicit construction.
\end{theorem}

\begin{proof}
The integral $\int g\, d\nu = \mathbb{E}_\nu[\log_2 3 - E] = \log_2 3 - \mathbb{E}_\nu[E]$.
By assumption $\mathbb{E}_\nu[E] = 2$
(from Proposition~\ref{prop:dist_E} in the i.i.d.\ model).
Thus the $\nu$-ergodic mean of $g$ is $\log_2 3 - 2 = \log_2(3/4) \approx -0.415 < 0$.
Since $\nu$ is $S$-invariant and ergodic, the Birkhoff ergodic theorem gives the stated
convergence for $\nu$-a.e.\ $x$.
\end{proof}

\begin{lemma}[Correspondence principle in the Collatz context]\label{lem:furstenberg}
Let $(X, \mathcal{B}, \mu, T)$ be a measure-preserving system and $A \in \mathcal{B}$
with $\mu(A) > 0$. A correspondence principle (cf.~\cite{Tao2022}) provides a structural result
connecting combinatorial properties of the Collatz orbit with the ergodic structure of
$(\mathbb{Z}_2, S, \mu)$.
In the Collatz context, the principle is used to transfer density results from the
2-adic setting to natural density statements about $\mathbb{N}$.
\end{lemma}

\section{Why Tao's Result Falls Short of the Conjecture}

\begin{remark}[The gap between logarithmic density and full conjecture]\label{rem:gap}
Tao's theorem establishes that the set of ``badly behaved'' orbits has logarithmic density 0.
However, the full Collatz conjecture requires showing that this set is \emph{empty} (or,
equivalently, that every orbit eventually reaches 1). These are very different statements:
a set of logarithmic density 0 can still be infinite (indeed, of natural density 0 but
infinite).
\end{remark}

\begin{proposition}[Obstacles to a full proof]\label{prop:obstacles}
The following obstacles prevent Tao's approach from proving the full conjecture:
\begin{enumerate}
  \item \textbf{Individual vs.\ statistical}: Tao controls the average behaviour but not
        the behaviour of any individual orbit.
  \item \textbf{No recurrence}: The 2-adic Fourier method gives density-zero results but
        not recurrence for every starting point.
  \item \textbf{Cycle exclusion}: Tao's argument does not rule out the existence of
        extremely large cycles that might ``trap'' some orbits.
  \item \textbf{The $+1$ barrier}: The $+1$ in $3n+1$ creates arithmetic correlations
        that destroy perfect independence of orbit steps, limiting the strength of the
        Fourier approach.
\end{enumerate}
\end{proposition}

\begin{conjecture}[Collatz conjecture, full version]\label{conj:full}
For every $n \in \mathbb{N}$, there exists $k \geq 0$ such that $T^k(n) = 1$.
\end{conjecture}

\section{Discussion and Open Problems}

Tao's result is significant in several ways:
\begin{enumerate}
  \item It establishes the strongest known quantitative statement toward the conjecture.
  \item It shows that if counterexamples exist, they must form a very ``sparse'' set
        (logarithmic density 0).
  \item It provides a toolkit (Fourier analysis over $\mathbb{Z}_2$, Syracuse random walks)
        that may be further developed.
\end{enumerate}

Open problems include:
\begin{enumerate}
  \item Can the logarithmic density result be improved to natural density?
  \item Can one show that divergent orbits have natural density 0 (a stronger statement)?
  \item Can the Fourier-analytic approach be combined with $p$-adic algebraic methods
        to exclude all non-trivial cycles?
  \item Is there a ``structural'' reason (beyond probabilistic heuristics) why the
        $+1$ forces all orbits to converge?
\end{enumerate}

\begin{thebibliography}{99}

\bibitem{Tao2022}
T.~Tao.
\newblock Almost all Collatz orbits attain almost bounded values.
\newblock \emph{Forum of Mathematics, Sigma}, 10:e72, 2022.

\bibitem{Collatz1937}
L.~Collatz.
\newblock Unpublished problem, communicated at the International Congress of Mathematicians, 1937.

\bibitem{Lagarias1985}
J.~C. Lagarias.
\newblock The $3x+1$ problem and its generalizations.
\newblock \emph{American Mathematical Monthly}, 92(1):3--23, 1985.

\bibitem{Lagarias2010}
J.~C. Lagarias (ed.).
\newblock \emph{The Ultimate Challenge: The $3x+1$ Problem}.
\newblock American Mathematical Society, Providence, RI, 2010.

\bibitem{Terras1976}
R.~Terras.
\newblock A stopping time problem on the positive integers.
\newblock \emph{Acta Arithmetica}, 30(3):241--252, 1976.

\bibitem{Koblitz1984}
N.~Koblitz.
\newblock \emph{$p$-adic Numbers, $p$-adic Analysis, and Zeta Functions}.
\newblock Springer-Verlag, New York, 1984.

\bibitem{Kontorovich2010}
A.~Kontorovich and J.~C. Lagarias.
\newblock Stochastic models for the $3x+1$ and $5x+1$ problems and related problems.
\newblock In \emph{The Ultimate Challenge}, pp.\ 131--188, AMS, 2010.

\end{thebibliography}

\end{document}
